% I. Анализ на предметната област и литературен преглед.
% Задача на раздела: да покаже къде стои проектът спрямо стандарта SQL, спрямо
% изградените върху него технологии за достъп и спрямо самите мрежови протоколи.
\section{Анализ на предметната област и литературен преглед}
\label{sec:literature}

\subsection{Стандартът SQL и мястото на клиентския интерфейс в него}

Стандартът ISO/IEC 9075 описва не само езика за заявки, но и начина, по който външна
програма стига до данните. Част 1 задава общата рамка и понятията, с които останалите
части си служат~\cite{sql_framework}. Част 3 определя \term{интерфейса на ниво
извикване} (Call-Level Interface, CLI), тоест процедурен интерфейс, чрез който
приложение на общоцелеви език изпраща заявления и получава резултати, без тези
заявления да са вградени в кода при компилация~\cite{sql_cli}. Именно тази част е
шаблонът, върху който стъпват практическите технологии за достъп.

Разграничението, което стандартът прави, е съществено за целия проект. Стандартът
нормира \emph{интерфейса}, който приложението вижда, но не нормира \emph{протокола},
по който клиентът и сървърът си говорят по мрежата. Протоколът е дело на всеки
доставчик поотделно. Следователно преносимостта, която CLI обещава, се постига винаги
чрез междинен слой, който превежда стандартните извиквания към нестандартния протокол
на конкретния сървър.

\subsection{Технологии за достъп, изградени върху стандарта}

ODBC е най-разпространената реализация на идеята за CLI. Приложението се свързва с
\term{диспечер на драйвери}, който зарежда драйвер за конкретния източник на данни и
пренасочва извикванията към него~\cite{odbc_ref}. Върху този модел исторически се
надграждат OLE DB, ADO и по-късно ADO.NET, които добавят обектен модел, работа с
несвързани множества от редове и по-богата типова система~\cite{adonet}. Общото между
всички тях е броят на границите, през които минава един резултатен ред, преди да
стигне до приложението.

Отделно от релационния път стоят спецификациите за слабо структурирани данни и за
директорийни услуги. XML остава нормативната основа за обмен на документно
структурирани данни между разнородни системи~\cite{xml10}, а LDAP описва достъпа до
йерархични директорийни хранилища с протокол, който по устройство е близък до
разглежданите тук: двоичен обмен на съобщения с ясно определени етикети и дължини,
върху постоянна връзка~\cite{rfc4511}. Сравнението с LDAP е полезно, защото показва,
че прякото говорене по протокол е обичайна практика извън релационния свят.

\subsection{Мрежовите протоколи на PostgreSQL и MySQL}

Протоколът на PostgreSQL е съобщително ориентиран: всяко съобщение носи еднобайтов
етикет, дължина и тяло, а фазите на работа са ясно отделени, стартиране и
удостоверяване, прост цикъл на заявка, разширен цикъл с
\code{Parse}/\code{Bind}/\code{Execute} за подготвени заявления~\cite{pg_protocol}.
Протоколът на MySQL е организиран около пакети с последователен номер и разчита в
по-голяма степен на позиционно кодиране и на цели числа с променлива
дължина~\cite{mysql_protocol}. Двата протокола решават една и съща задача с различни
средства, което ги прави удобна двойка за общ клиентски слой.

Референтните клиентски библиотеки на двата доставчика са добре документирани и служат
за еталон в тази работа. Документацията на \code{libpq} описва подробно както
интерфейса, така и семантиката на асинхронния режим~\cite{libpq}. По-общата литература
за системи за обработка на транзакции обяснява защо клиентският слой не е неутрален
към производителността: границите на транзакциите, изчакванията и броят обръщания по
мрежата определят наблюдаваната закъснителност повече от скоростта на самия
сървър~\cite{gray1993, kleppmann2017}.

\subsection{Изводи от прегледа}

Прегледът очертава три нива, на които може да се стъпи: стандартизиран интерфейс с
диспечер на драйвери, доставчикова клиентска библиотека и самата спецификация на
протокола. Всяко следващо ниво надолу премахва слой и абстракция, но добавя
отговорност за коректност. \TODO{да се добавят два или три съпоставими академични или
индустриални източника за измервана цена на слоевете при достъп до бази данни, ако
такива бъдат намерени; ако не бъдат, това да се заяви изрично тук}
